
\section*{7. Generated subgroups; commutator group}
\phantomsection\label{sec:sga3-VIB-7}%
\addcontentsline{toc}{section}{7. Generated subgroups; commutator group}

Throughout this section, \(k\) denotes a fixed field.

\medskip
\phantomsection\label{prop:sga3-VIB-7.1}%
\noindent\textbf{Proposition 7.1.---}
\emph{Let \(G\) be a \(k\)-group, let
\((X_i)_{i\in I}\) be a family of geometrically reduced
\(k\)-schemes,}%
\emph{and, for every \(i\in I\), let \(f_i:X_i\to G\) be a
\(k\)-morphism.}
\begin{enumerate}
\item[\textup{\SGAInlineRefTarget{prop:sga3-VIB-7.1:item:i}(i)}]\phantomsection\label{concl:sga3-VIB-7.1-i}
  \emph{There exists a smallest closed \(k\)-group subscheme of
  \(G\) through which every \(f_i\) factors. It is denoted
  \[
    \Gamma_G\bigl((f_i)_{i\in I}\bigr).
  \]
  It is a geometrically reduced \(k\)-scheme, and hence smooth
  when \(G\) is locally of finite type over \(k\)
  (\hyperref[cor:sga3-VIB-1.3.1]{1.3.1}).}
\item[\textup{\SGAInlineRefTarget{prop:sga3-VIB-7.1:item:ii}(ii)}]\phantomsection\label{concl:sga3-VIB-7.1-ii}
  \emph{Put
  \[
    X=\coprod_{i\in I}X_i,
  \]
  and let \(f:X\to G\) be the morphism whose restriction to
  \(X_i\) is \(f_i\), for every \(i\in I\). Put
  \[
    X^1=X\amalg X,
  \]
  and let \(f^1:X^1\to G\) be the morphism whose restrictions to
  the two copies of \(X\) are \(f\) and \(c\circ f\),
  respectively. For every \(n>1\), put
  \[
    X^n=X^1\times_kX^{n-1},
    \qquad
    f^n=\mu\circ(f^1\times_k f^{n-1}):X^n\longrightarrow G.
  \]
  Then \(\Gamma_G((f_i)_{i\in I})\) is the reduced subscheme of
  \(G\) whose underlying space is the closure of
  \[
    \bigcup_{n\geq1}f^n(X^n).
  \]}
\item[\textup{\SGAInlineRefTarget{prop:sga3-VIB-7.1:item:iii}(iii)}]\phantomsection\label{concl:sga3-VIB-7.1-iii}
  \emph{For every \(k\)-scheme \(S\),
  \(\Gamma_G((f_i)_{i\in I})_S\) is the smallest closed group
  subscheme of \(G_S\) through which every
  \[
    f_{i,S}:(X_i)_S\longrightarrow G_S
  \]
  factors.}
\item[\textup{(iii\(^{\prime}\))}]
  \phantomsection\label{concl:sga3-VIB-7.1-iiiprime}
  \emph{Moreover, \(\Gamma_G((f_i)_{i\in I})_S\) is the smallest
  group subscheme of \(G_S\), not necessarily closed, through
  which every \(f_{i,S}\) factors.}%
\end{enumerate}

First note that, to prove \textup{\SGARefLink{prop:sga3-VIB-7.1:item:i}{(i)}}, \textup{\SGARefLink{prop:sga3-VIB-7.1:item:iii}{(iii)}}, and
\textup{(iii\(^{\prime}\))}, after defining \(X\) and \(f\) as
in \textup{\SGARefLink{concl:sga3-VIB-7.1-ii}{(ii)}}, we are reduced to the case in which \(I\)
consists of a single element.

Let \(H\) be the reduced subscheme of \(G\) whose underlying set
is
\[
  \overline{\bigcup_{n\geq1}f^n(X^n)}.
\]
The family of morphisms
\[
  f^n:X^n\longrightarrow H
\]
is schematically dominant (cf.\ EGA~IV\(_3\), 11.10.4), and therefore
every closed subscheme of \(G\) through which all the \(f^n\)
factor also contains \(H\). Moreover, by \emph{loc.\ cit.},
11.10.7, \(H\) is geometrically reduced. To prove
\textup{\SGARefLink{prop:sga3-VIB-7.1:item:i}{(i)}} and \textup{\SGARefLink{prop:sga3-VIB-7.1:item:ii}{(ii)}}, it is therefore enough to show
that \(H\) is a group subscheme of \(G\), that is, that the
restriction of \(c\) to \(H\) and the restriction of \(\mu\) to
\(H\times_kH\) factor through the inclusion \(H\to G\).

Since \(H\) is geometrically reduced, \(H\times_kH\) is reduced,
and it is enough to verify, set-theoretically, that
\[
  c(H)\subset H,
  \qquad
  \mu(H\times_kH)\subset H.
\]
By EGA~IV\(_3\), 11.10.6, the union of the
\[
  f^n_{(H)}(X^n\times_kH)
\]
is schematically dense in \(H\times_kH\). Similarly, for every
\(n\geq1\), the union of the
\[
  f^m_{(X^n)}(X^n\times_kX^m),
  \qquad m\geq1,
\]
is schematically dense in \(X^n\times_kH\). It is therefore
enough to prove
\[
  \mu\!\left(
    f^n_{(H)}
      \bigl(f^m_{(X^n)}(X^n\times_kX^m)\bigr)
  \right)\subset H
\]
and \(c(f^n(X^n))\subset H\). But
\[
\begin{aligned}
  \mu\!\left(
    f^n_{(H)}
      \bigl(f^m_{(X^n)}(X^n\times_kX^m)\bigr)
  \right)
  &=
  \mu\!\left((f^n\times_k f^m)(X^n\times_kX^m)\right)\\
  &=f^{n+m}(X^{n+m})\subset H.
\end{aligned}
\]
Also, since \(c(f^1(X^1))\subset f^1(X^1)\), for every \(n\)
one has
\[
  c(f^n(X^n))\subset f^n(X^n)\subset H.
\]
This proves \textup{\SGARefLink{prop:sga3-VIB-7.1:item:i}{(i)}} and \textup{\SGARefLink{prop:sga3-VIB-7.1:item:ii}{(ii)}}.

We now prove \textup{\SGARefLink{concl:sga3-VIB-7.1-iii}{(iii)}}. Let \(G'\) be a closed group
subscheme of \(G_S\) through which \(f_S\) factors. We must show
that \(G'\) contains \(H_S\), or equivalently that
\[
  H_S=H_S\times_{G_S}G'.
\]
Put
\[
  H'_S=H_S\times_{G_S}G'=G'\cap H_S.
\]
Every \(f^n_S\) factors through both \(G'\) and \(H_S\), and
hence through \(H'_S\). Since the family
\[
  f^n:X^n\longrightarrow H
\]
is schematically dominant, EGA~IV\(_3\), 11.10.6 implies that
the base-changed family
\[
  f^n_S:X^n_S\longrightarrow H_S
\]
is schematically dominant as well. Because every \(f^n_S\)
factors through \(H'_S\), EGA~IV\(_3\), 11.10.1\textup{(c)}
gives
\[
  H'_S=H_S.
\]
This proves \textup{\SGARefLink{prop:sga3-VIB-7.1:item:iii}{(iii)}}.

Finally, we show that \(H_S\) is the smallest group subscheme,
not necessarily closed, of \(G_S\) through which \(f_S\)
factors. Let \(G'\) be such a group subscheme and again put
\[
  H'_S=H_S\times_{G_S}G'.
\]
We must show \(H'_S=H_S\). By \textup{\SGARefLink{concl:sga3-VIB-7.1-iii}{(iii)}}, it is enough to
show that \(H'_S\) is closed in \(H_S\). It is enough in turn to
show that \(H_S\) and \(H'_S\) have the same underlying set, and
for this it suffices to show, for every \(s\in S\), that
\[
  H'_s
  :=H'_S\times_S\kappa(s)
  =H_s\times_{G_s}G'_s
  =H_s.
\]
By \hyperref[sga3:VIA:0.5.2]{VI\(_A\), 0.5.2}, the
\(\kappa(s)\)-group subscheme \(G'_s\) is closed in \(G_s\).
Hence \(H'_s\) is closed in \(H_s\), and the preceding argument,
applied to \(H'_s\), \(H_s\), and the \(f^n_s\), shows that
\[
  H'_s=H_s.
\]
This proves \textup{(iii\(^{\prime}\))} and completes the proof.
